\section{Preprocessing}

The preprocessing pipeline converts raw text corpora into fixed-length token
chunks suitable for autoregressive language model training.  It is designed to
support multi-stage curriculum learning with replay buffers and
difficulty-aware sampling.

\subsection{Tokenization and Chunking}

Raw documents are streamed from HuggingFace datasets and processed through a
BPE tokenizer (Section~\ref{sec:tokenizer}).  Each document is wrapped with
\texttt{BOS} (id=2) and \texttt{EOS} (id=3) tokens, producing the sequence:
\[
[\text{BOS}] \;\|\; t_1, t_2, \ldots, t_n \;\|\; [\text{EOS}]
\]
These per-document token sequences are concatenated into a single stream, then
sliced into non-overlapping windows of length $(\texttt{seq\_len} + 1)$.  The
extra token provides the target for the final position (input--target shift).
The last partial window is discarded.

\textbf{Sequence lengths per stage:} Stage~0--3 use $\texttt{seq\_len}=384$,
Stages~4--5 use 512, and Stage~6 uses 768, matching the model's context length
at each expansion.

\subsection{Caching}

All chunked datasets are serialized to disk as \texttt{.pkl} files keyed by
\texttt{(dataset\_name, seq\_len)}.  Subsequent runs load directly from cache,
avoiding re-tokenization.  Mixed-dataset caches use an MD5 hash of the
sorted dataset--weight dictionary as part of the filename to ensure
deterministic cache keys.

\subsection{Difficulty Scoring}

For curriculum-aware modes (\texttt{perplexity}, \texttt{adaptive}), each
story is pre-scored using a composite difficulty metric computed by a
reference model.  Scores are saved in \texttt{curriculum\_scores.npy} as a
$(N, 5)$ array.  Chunk-level difficulties are derived by interpolating story
scores to chunk positions based on their sequential ordering in the dataset.

\subsection{Curriculum Filtering}

In \texttt{adaptive} mode, only a fraction of chunks (sorted by difficulty) is
exposed to the model at any time.  The \texttt{CompetenceScheduler} adjusts
this fraction based on validation improvement rates:
\begin{itemize}[nosep]
  \item Strong improvement ($>$1\% loss decrease) $\to$ expand pool by +5\%
  \item Moderate improvement ($>$0.1\%) $\to$ +2\%
  \item Stagnation for 3 consecutive checks $\to$ forced +1\%
\end{itemize}
The expansion stages use \texttt{random} mode (no curriculum sorting), since
the difficulty gradient is already encoded in the staged domain progression.

\subsection{Anchor Injection}

Regardless of curriculum mode, 10\% of training samples are drawn from an
\emph{anchor pool} consisting of the easiest 5\% of chunks.  This provides
continuous exposure to simple patterns and helps stabilize gradient norms.

\subsection{Replay Pool Preparation}

Each stage loads pre-cached chunks from earlier corpora as a replay buffer.
Source chunks are loaded from the largest available cache file, then:
\begin{itemize}[nosep]
  \item Chunks longer than the target \texttt{seq\_len} are truncated.
  \item Chunks shorter than the target are skipped (not zero-padded, to avoid
        corrupting the loss signal).
\end{itemize}
Replay samples are drawn with probability \texttt{replay\_frac} at each batch
construction, replacing a training chunk with one from the replay pool.
